% Nouveau EGA V (cote 25): archive cote 25, site page 29 (PDF page 30).
% typescript. Draft compared: none.
% AI-transcribed by Claude Opus 5.5 (claude-opus-5-5, adaptive thinking, effort high), Message Batches API, batch ega25b (msgbatch_01A7n5YuXubUPNK1zsepi8RG), 2026-11-13; not yet checked by eye.
\scanpage{25/29}
\mspage{\hand{25}}

Je dois cependant d'abord \unsure{signaler|signqler} un remords à 4.2. et 4.3. , qui
\unsure{dev-aient|devaient} être énoncés sous une forme légèrement plus générale (tout au
moins en remarque): on se donne un entier m, et on \unsure{désigne|dásigne} par E l'ensem
ble des \hand{$\xi\in$}$\overset{\hand{$\vee$}}{\text{P}}$ tels que G$_{\hand{$\xi$}}$ resp. Y$_{\hand{$\xi$}}$ satisfasse P sauf \deleted{au plus} sur une
partie fermée de dimension\hand{$\leq$}m (i.e. l'ensemble P-singulier Z est de dimen-%
sion\hand{$\leq$} m). Alors a) E est une partie constructible de $\overset{\hand{$\vee$}}{\text{P}}$ et b) dans le
cas S=Spec(k), si F \inserted[typed]{resp. X} satisfait à P sauf \deleted{au plus} sur un ensemble de
\xout{xx}dimension\hand{$\leq$} m+1 , alors E contient un ouvert non vide.

\uline{Proposition} 7.2. Supposons que X soit propre sur S, et F de présentati-%
on finie et plat sur S. Soit P l'une des propriétés (i) à (v) de 4.2.
et soit m un entier. Soit s\hand{$\in$}S, et supposons que \xout{k(s)xxF$_s$xsatisfassexà}
\xout{laxpropriétéxP} l\xout{'}ensemble Z$_s$ des points de X$_s$ en lesquels F$_s$ ne satisfait
pas à P soit de dimension \hand{$\leq$} m+1. Alors, si de plus k(s) est infini, il
existe un voisinage U de s dans E, et une section \hand{$\xi$} de $\overset{\hand{$\vee$}}{\text{P}}$ sur U, ayant
les propriétés suivantes: pour tout s\hand{$\in$}U, l'ensemble Z des points
de \deleted{X} Y$_{\hand{\unsure{$\xi,\xi(s)$|$\xi(s)$}}}$ en lesquels G$_{\hand{$\xi(s)$}}$ ne satisfait pas à P est de dimension
\hand{$\leq$}m, et $\phi_{\hand{$\xi$}(s)}$ est F$_{\hand{\unsure{$\xi,\xi(s)$|$\xi(s)$}}}$-régulière\hand{,} \xout{xlexModulexGxxsur} Sous ces condi-%
tions, le Module G$_{\hand{$\xi$}}$ sur Y$_{\hand{$\xi$}}$ est plat relativement à U. Enfin, si k(s)
n'est pas supposé \unsure{inffni|infini}, on peut encore faire la construction précéden-%
te, après extension étale duntype envisagé dans 7.1.

\uline{Proposition} 7.3. Essentiellement la même, il n'y a plus de F, \unsure{on|ok}
\margin{[N.B. Pour (v) cf 12.2.1. (x) et (xi) \unsure{Ceci peut} donc apporter une variante à 4.3. \unsure{via} 12.2.1. (x)]}
suppose X plat rel à S, on réfère aux propriétés (i) à \deleted{(xxx)}\inserted{\unsure{(iv)}} de 4.3. au
lieu de celles de 4.2. (mais en ayant soin de faire \deleted{les réserves} \xout{axxxxxxxxxxx}\inserted{la réserve}
``k\inserted{(s)} de car. 0 ou f:X$\longrightarrow$P une immersion\hand{$\xi$}\inserted{, et dans le cas (v), pour tout $s\in S$, et tte comp. irréductible Z de X$_s$, on ait \unsure{dif}(Z) $\geq 2$).}

\deleted{On notera que nous avons exclus les propriétés (v) et (vi) de}
\deleted{4.3. , qui ne sont pas de nature purement locale. Cependant, on récupère}
\deleted{le cas (v) en notant qu'on peut s'arranger en vertu de la démonstratio}
\deleted{de 7.3. que Y$_{\hand{$\xi$}}$ soit plat sur U à fibres séparables, mais alors en}
\margin{\unsure{29}}
